\documentclass[10pt,a4paper]{amsart}
\usepackage[margin=19mm]{geometry}
\usepackage{amsmath,amssymb,mathtools}
\usepackage[scr=rsfs,cal=euler]{mathalfa}
\usepackage{xfrac}
\usepackage[unicode,hidelinks,bookmarks=false]{hyperref}
\newcommand{\ie}{i.e.\ }
\newcommand{\cf}{cf.\ }
\newcommand{\resp}{respectively\ }
\newcommand{\Subsection}{\subsection}
\providecommand{\nobreakdash}{\nobreak\hbox{-}}
\setlength{\emergencystretch}{2em}
\multlinegap0pt
\usepackage{amssymb}
\usepackage{bm}
\usepackage{mathtools}
\usepackage{tikz}
\usepackage[shortlabels]{enumitem}
\usepackage{algorithm}\usepackage{algpseudocode}
\newtheorem{lem}{Lemma}
\newcommand{\C}{\SET{C}}
\newcommand{\QU}[1]{\mathcal{#1}}
\newcommand{\SET}[1]{\mathbb{#1}}
\newcommand{\pure}[2]{\C^{#1 \otimes #2}}
\newcommand{\setint}[1]{[\hspace{-0.5mm}[#1]\hspace{-0.5mm}]}
\newcommand{\state}[2]{S(\pure{#1}{#2})}
\newcommand{\subint}[2]{\tbinom{\setint{#1}}{#2}}
\title{Insertion Correcting Capability for Quantum Deletion-Correcting Codes}
\author{Ken Nakamura, Takayuki Nozaki}
\date{Source: arXiv:2602.20635v2 (CC BY 4.0)}
\begin{document}
\maketitle

\noindent\emph{Source and licence.} Insertion Correcting Capability for Quantum Deletion-Correcting Codes. Version arXiv:2602.20635v2 (CC BY 4.0), distributed by arXiv under the Creative Commons Attribution 4.0 International licence, http://creativecommons.org/licenses/by/4.0/, as recorded in the arXiv licence metadata for this exact version and shown on the abstract page https://arxiv.org/abs/2602.20635v2. Full licence notice: \texttt{licenses/arxiv-2602.20635-CC-BY-4.0.txt}.\par\smallskip

\noindent\textit{Editorial scope.} Counted unit: the article's lem lem:1ins\_1del\_kai (source0.tex lines 1066--1074). The article's complete original proof of it is delivered with it (source0.tex lines 1230--1286, one \texttt{proof} environment of 57 lines, which the article places away from the statement). No mathematics is written, repaired or re-derived: the statement and the proof are the article's own lines, in the article's own notation, and no retained line is substituted or re-flowed.  Retained uncounted as the source-local context the proof needs: lines 1201--1210.  Nothing was omitted from any retained span, so the package carries no omission record.  No retained line contains a citation, so the wrapper carries no bibliography and no bibliography entry.  No retained line contains a figure environment, an image-inclusion command or drawing code, so no image is delivered and the package declares no asset dependency.  Editorial formatting only, in addition: an AMS \texttt{amsart} preamble that restates the article's own declarations and macros used by the retained lines, namely the environments \texttt{lem} on separate counters, together with the standard packages those lines need.  The retained environments print in this document as Lemma 1 in order of appearance, whereas the article prints the same items with the numbering of its own full document.  The complete native source of the article as distributed in the arXiv e-print for this exact version is bundled unchanged as \texttt{sources/195204/source0.tex}.
\begin{lem}
  \label{lem_ID=non}
  For $P \in \subint{n}{t}$,
  $\QU{X} \subset \state{l}{n}$ satisfies
  \begin{align*}
    \QU{X}
    \subset
    \QU{I}_P \circ \QU{D}_P(\QU{X}).
  \end{align*} 
\end{lem}

\begin{lem}
  \label{lem:1ins_1del_kai}
  Any set $\QU{X} \subset \state{l}{n}$ satisfies
  \begin{align*}
    \QU{D}^1 \circ \QU{I}^1(\QU{X})
    \subset
    \QU{I}^1 \circ \QU{D}^1(\QU{X}).
  \end{align*}
\end{lem}

\begin{proof}[Proof of Lemma \ref{lem:1ins_1del_kai}]
  To simplify the notation,
  we write $\QU{D}_{\{ p \}}$ as $\QU{D}_p$
  and
  $\QU{I}_{\{ q \}}$ as $\QU{I}_q$.
  It suffices to prove that
  for $\rho \in \state{l}{n}$,
  \begin{align}
    \forall
    p,q \in \setint{n+1}
    ~~~
    \QU{D}_q \circ \QU{I}_p( \rho )
    \subset
    \QU{I}^1 \circ \QU{D}^1( \rho ).
    \label{eq:juubunn}    
  \end{align}
  
  If $n = 1$,
  Eq.~\eqref{eq:juubunn} holds
  since
  $\QU{I}^1 \circ \QU{D}^1( \rho ) = \state{l}{1}$.
  Hence,
  consider the case for $n \ge 2$.
  If $p = q$,
  from Lemma \ref{lem_ID=non},
  Eq.~\eqref{eq:juubunn} follows
  \begin{align*}
    \QU{D}_p \circ \QU{I}_p( \rho )
    &=
    \{ \rho \}
    \subset
    \QU{I}_p \circ \QU{D}_p( \rho ). 
  \end{align*}
  For $p < q$,
  by substituting $\QU{X} = \QU{D}_q \circ \QU{I}_p( \rho )$ into Lemma \ref{lem_ID=non},
  we have 
  \begin{align*}
    \QU{D}_q \circ \QU{I}_p( \rho )
    &\subset
    \QU{I}_p \circ \QU{D}_p \circ \QU{D}_q \circ \QU{I}_p( \rho )\\
    &=
    \QU{I}_p \circ \QU{D}_{q-1} \circ \QU{D}_p \circ \QU{I}_p( \rho )\\
    &=
    \QU{I}_p \circ \QU{D}_{q-1}( \rho )\\
    &\subset
    \QU{I}^1 \circ \QU{D}^1( \rho ),
  \end{align*}
  where
  the first equality
  follows
  from $\QU{D}_p \circ \QU{D}_q = \QU{D}_{q-1} \circ \QU{D}_p$.
  Similarly,
  we get Eq.~\eqref{eq:juubunn}
  for $p > q$.
  Thus,
  we get Eq.~\eqref{eq:juubunn}.
\end{proof}

\end{document}
